% ===== 4. Methods =====
\section{4. Methods}\label{methods}

\textbf{Subject model and generation.} The subject is Qwen3-1.7B
(\texttt{Qwen/Qwen3-1.7B}; \citealp{qwen2025qwen3}), a 1.72B-parameter
text-only instruction-tuned decoder-only transformer (Apache 2.0), with
28 blocks and hidden dimension 2048. I run it in bfloat16 on Apple
Silicon MPS in 16 GB unified memory. Decoding is greedy throughout
(\texttt{do\_sample=False}, \texttt{max\_new\_tokens=64},
\texttt{enable\_thinking=False}, suppressing Qwen3's
\texttt{\textless{}think\textgreater{}} trace), so all generations are
deterministic conditional on the prompt.

\textbf{Hidden-state capture.} Each question runs
\texttt{model.generate(...)} with \texttt{output\_hidden\_states=True}
and \texttt{output\_scores=True}, persisting one \((L+1, d) =
(29, 2048)\) prefill tensor per question. One indexing fact matters
downstream. Entries \(1..L{-}1\) of the returned stack are raw
post-block residuals, but the final entry (index \(L = 28\)) is the
\emph{final normed} state, because HF applies the last RMSNorm before
recording it. So the probes read the normed state at ``layer 28'' while
the Section 6.2/6.3.1 hooks and the SAE operate on the pre-norm
post-block-27 residual.

\textbf{Same-model judge.} Qwen3-1.7B labels each generated answer
itself, called separately with a structured prompt (\texttt{judge.py})
that asks for one of \texttt{correct}, \texttt{refusal}, or
\texttt{wrong} (decision rule, worked examples, and parser in Appendix
B.1). This judge replaces the brittle substring-match grader used in
earlier work.

Three LLM annotators measure judge agreement on the full benchmark,
applying the same three-way rule strictly (Table \ref{tbl:judge},
Appendix B.3): the Qwen3-1.7B judge, Claude Opus 4.7 on a stratified
\(30\)-item sample, and Google Deep Research (DR; Gemini 3.1 Pro with
web-search verification) covering all \(784\). DR vs.~the Qwen-judge
gives \(\kappa = 0.78\), ``substantial agreement''
\citep{landis1977agreement}. Claude agrees \(30/30\), which bounds the
Qwen-judge's disagreement with a second independent grader at
\(\le 10\%\) at \(95\%\) confidence by the rule of three. DR is
systematically stricter, reflecting its containment rule and gold
verification rather than judge miscalibration (decomposition and
retained flagged-premise items in Appendix B.3). A sensitivity rerun
under the DR labels keeps the picture, and a blind Claude-family regrade
of the Section 6 intervened generations agrees with the self-judge
without inflating flip rates (both in Appendix B.3).

\textbf{Hidden-state probe pipeline.} For each binary target and layer,
I fit the Section 2.3 pipeline on the last-prompt-token states at that
depth. I report mean and std of fold accuracy plus a ``layer range'',
the set of layers within \(1\sigma\) of the peak.

\textbf{Prompt-feature baseline pipeline.} For each target, I also fit
baselines with no access to the hidden state, reporting 5-fold CV
accuracy on the \emph{same folds} (\texttt{random\_state=0}). Four
baselines of increasing leniency: \textbf{TF-IDF} (bag-of-tokens
logistic regression on the raw question text); \textbf{text-only}
(engineered numeric features); \textbf{text + domain} (adds the domain
one-hot); and \textbf{text + domain + category} (adds the
\emph{annotator's} expected-difficulty dummy, the least conservative
baseline). For the \texttt{cutoff} target the category dummy is excluded
from the strongest-baseline comparison, since it would be the label
itself. Full feature lists and hyperparameters are in Appendix B.7. The
honest comparison is the hidden-state probe vs.~the \emph{strongest} of
the four on each target: the ``h adds'' column of Table \ref{tbl:hadds} in Appendix A.5.

\textbf{Probe targets.} \texttt{correct} (all 784 items);
\texttt{cutoff} (post vs.~pre, all 784); \texttt{refusal\_vs\_wrong}
(the \texttt{\{refusal,\ wrong\}} subset, n=549);
\texttt{refusal\_vs\_wrong\_within\_post} (n=393; all 147 refusals are
post-cutoff, so this conditions out the cutoff covariate);
\texttt{correct\_within\_pre} (n=296, the cutoff disconfound test); and
\texttt{correct\_within\_obscure} (n=153, the popularity disconfound
test).

\textbf{Bootstrap protocol (Sections 8.1, 8.2).} For each
\texttt{(dataset,\ model,\ target)} cell, I draw \(K = 30\) random 50/50
class-balanced subsamples without replacement. Each refits the full
probe pipeline (peak over all \(29\) layers) plus the four prompt
baselines on the same folds, and I report \(\bar h_{\rm adds}\) with a
percentile 95\% CI (class-size rule, CI construction, and compute bounds
in Appendix E). One asymmetry matters when reading small margins: the
probe side maximizes over \(29\) correlated layers, the baseline side
over four. Under a no-signal null the expected \(h_{\rm adds}\) is
therefore slightly positive, so margins within a couple of points of
zero carry this selection bias.

\textbf{External-dataset evaluation (Section 8.1).} \textbf{PopQA}
\citep{mallen2023popqa} and \textbf{TriviaQA} \citep{joshi2017triviaqa} run
end-to-end through the same generation pipeline and three-way
Qwen3-1.7B judge, \(n=800\) per seed (splits, strata, seeds in Appendix
E). Three seeds per \texttt{(subject\ model,\ dataset)} widen the pool
to \(n_{\text{unique}} \approx 2{,}200\); Llama runs are single-seed
(\(n_{\text{unique}} = 800\)). The judge remains Qwen3-1.7B in all runs
(Section 10 records the Qwen-judge-on-Llama-outputs caveat).
